\section{Discussion and Limitations}
\label{sec:discussion}

SurgAlign is best understood as a supervision-refinement method. It targets one specific failure mode of weak video-language pretraining: a video-text pair can be semantically valid even when only a small subset of frames provides the visual evidence that justifies it. For that reason, the method is intentionally training-only. We do not introduce a new test-time retrieval head or a query-conditioned inference pipeline; the claim is simply that weak temporal positives should be handled more carefully during pretraining.

The evaluation protocol is designed around the first limitation. SurgAlign is most appropriate when the paired text remains semantically informative, the temporal offset is local rather than arbitrarily long-range, and the decisive evidence is sparse enough that uniform pooling becomes harmful. The method is therefore not a substitute for stronger backbones, richer negatives, or better corpus design.

The second limitation is empirical scope. Our main claim is tied to internal matched comparisons, with SurgCLIP-$\beta$ and full SurgCLIP treated as reported strong references rather than as matched baselines. This distinction matters because temporal misalignment is real but not exclusive: tasks dominated by broader workflow context, harder negative discrimination, or label ambiguity can require additional modeling beyond local frame weighting.

A third limitation is optimization sensitivity. In practice, the final behavior can depend on the expansion ratio, the auxiliary loss weight, and the training trajectory. To avoid converting this sensitivity into hidden test-set checkpoint selection, we fix the official checkpoint at epoch 50 and report multi-seed aggregates. Additional checkpoint trajectories are used only as diagnostic analyses.

A final limitation is evaluation breadth. We validate SurgAlign on surgical-domain benchmarks, where temporal misalignment is practically important and visually subtle, but we do not yet show transfer to broader narrated video datasets. The broader claim of the paper is therefore not that SurgAlign solves weak video-language alignment in general, but that treating weak intervals as noisy temporal anchors is a useful training principle in a challenging real-world regime.
